\section{Motivation}

\begin{frame}{Motivation for Self-Supervised Learning}
    \begin{enumerate}
        \setlength{\itemsep}{-0.5em}
        \item \textbf{Imagine you have tons of data, but very few labels.}\\
        Annotating data is expensive and slow, but unlabeled data is everywhere!\\[0.5em]
        \item<2-> \textbf{What if you could learn useful features from all that unlabeled data?}\\
        Self-supervised learning lets us pre-train models to extract general, reusable representations, reducing the need for labeled data later.\\[0.5em]
        \item<3-> \textbf{Think beyond images:}\\
        These techniques work not just for vision, but also for language, audio, and time-series data, enabling unified approaches across domains.\\[0.5em]
        \item<4-> \textbf{Inspired by how humans learn:}\\
        We learn patterns from the world around us without explicit labels; self-supervised learning mimics this natural process.\\[0.5em]
        \item<5-> \textbf{A journey of methods:}\\
        The field has evolved from generative models (like VAEs and GANs), to contrastive and predictive coding, and now to hybrid approaches combining reconstruction and discrimination.
    \end{enumerate}
\end{frame}

\begin{frame}[allowframebreaks]{}
    \begin{figure}
        \centering
        \fetchconvertimage{https://framerusercontent.com/images/r0i6jmANAR5bsgzvqm0BJG9aas.png}{images/ssl/architecture.png}{height=0.82\textheight,width=1\textwidth,keepaspectratio}
        \caption*{\scriptsize [Source: \url{https://framerusercontent.com/images/r0i6jmANAR5bsgzvqm0BJG9aas.png}]}
    \end{figure}
\end{frame}